To integrate formal libraries from various sources and to implement generic knowledge management services, as is one of the objectives of this thesis (see \autoref{sec:intro:objectives}), it is necessary to have a unifying framework and language in which these libraries are represented. Naturally, \ommt is (due to its generality and modularity) the prime candidate for such a framework. This Part covers the process of translating one library from an external system into the \ommt language.

\paragraph{} The OAF project (see \autoref{sec:intro:oaf}) in particular focuses on integrating libraries from theorem prover systems; however, the methodology involved generalizes easily to the purposes of the \ODK project (see \autoref{sec:intro:prelim:odk}), where the aim is to integrate databases of mathematical objects and their properties (e.g. finite groups), computer algebra systems and related systems. Together, they cover three aspects of the \emph{tetrapod} (see \autoref{ch:intro}), namely \emph{inference}, \emph{tabulation} and \emph{computation}. Additionally, the \mmt system itself manages the \emph{organization} aspect.

In all settings, the process involved in translating a library $L$ from system $S$ to \ommt entails the following steps:
\begin{enumerate}
	\item Use a suitable logical framework (as discussed in \autoref{ch:lf}) to formalize the foundational ontology of $S$.
	\item Obtain an export of the contents of $L$ in a parsable format. In the case of theorem prover libraries, this usually requires collaboration with the developers of $S$, since the source files themselves are often only readable for the corresponding system and do not contain inferred information such as implicit arguments, resolved notations, relevant type-check conditions, etc.
	\item Implement an importer that systematically translates the export from the previous step into \ommt, using the symbols in the formalization from step 1 to represent the foundational concepts of $S$.
\end{enumerate}

There is a continuum of possible translations, ranging from perfect adequacy to mere representation. In the best case, we can fully type check the result of the translation and the precise semantics of the original content is preserved under translation. In the worst case, we obtain a representation of the symbols in the original library as untyped and undefined constants. Ideally we want the former, however, depending on the specifics of the system, this might be impossible to achieve without reimplementing the entire system within \mmt. For example, a system with predicate subtypes and a powerful automated prover will usually rely on the latter to type check contents in a manner that can only be reproduced with a similarly powerful prover. In the case of (untyped) computer algebra systems, it might be impossible to check whether a given expression in the language of the system is even well-formed without passing it back to the system itself and evaluating the result.

However, it should be noted that for many knowledge management services such as alignments (see \autoref{sec:crosslib:alignments}) or (in a limited form) content translation (see \autoref{sec:crosslib:translate}), even mere representation is sufficient to obtain useful results. From that point of view, an adequate translation is not strictly necessary, but usually the usefulness of generic services strongly correlates with the amount of details and structure preserved in the translated library.

\subparagraph{} In this thesis we only consider the \emph{statements} in a given formal library (declarations, definitions, theorems, axioms) -- i.e. we do not cover \emph{proofs} or specific implementations of methods in computer algebra system. Importing proofs poses additional challenges and requires an adequate proof language sufficiently generic to subsume the various ways proofs can be represented in formal systems, if they can be exported from the system at all. To the best of my knowledge, such a generic proof language that can serve as an adequate target for translations from theorem prover systems does not yet exist. All knowledge management services presented in this thesis consequently do not benefit from any knowledge about proofs, other than (potentially) the binary information whether a statement is proven at all and (possiby) dependency management.

\paragraph{} The methodology described in this part has been applied to various systems and libraries by several people, resulting in \ommt imports for Mizar~\cite{IKRU:mizar:11}, HOL Light~\cite{KR:hollight:14}, TPTP~\cite{tptp}, IMPS~\cite{jbetzendahl:msc18}, Isabelle (as-of-yet unpublished), Coq~\cite{coqmmt} and PVS~\cite{KohMueOwr:mpagsiuf17}. As a representative example, the latter is covered in detail in \autoref{ch:import:pvs}. Moving on to less formal corpora, \autoref{sec:import:lmfdb} demonstrates how to handle databases of mathematical objects, exemplary the LMFDB, and \autoref{sec:import:cas} covers computer algebra systems, using the \GAP and \Sage systems as examples.